section{Distributed Ledger Technology}

\subsection{Definisi}
Secara fundamental, \textit{Distributed Ledger Technology} (DLT) adalah teknologi basis data terdesentralisasi yang memungkinkan transaksi terjadi secara \textit{peer-to-peer} tanpa memerlukan otoritas pusat atau perantara untuk memvalidasinya. Sistem ini dirancang secara spesifik untuk dapat beroperasi dalam "lingkungan adversarial", yaitu kondisi di mana partisipan jaringan mungkin tidak saling percaya atau bahkan berniat jahat, namun tetap mampu mencapai kesepakatan (konsensus) mengenai validitas data bersama \cite{16}.

Alih-alih bergantung pada pihak ketiga tepercaya, DLT mendistribusikan kepercayaan ke seluruh jaringan melalui mekanisme kriptografi, yang menjamin bahwa catatan yang telah divalidasi bersifat \textit{immutable} (tidak dapat diubah) serta meningkatkan efisiensi dan keamanan data. Dalam kerangka kerjanya, DLT berfungsi sebagai "mesin konsensus" yang memastikan integritas data tetap terjaga di seluruh salinan \textit{ledger} yang tersebar, meskipun tanpa adanya koordinator pusat. Penting juga untuk dipahami bahwa DLT adalah istilah payung yang lebih luas, di mana Blockchain hanyalah salah satu jenis atau subset dari teknologi ini yang menggunakan struktur data rantai blok \cite{16}.